\chapter*{\center \Large ANEXOS}
\addcontentsline{toc}{section}{\bfseries ANEXOS}
\markboth{ANEXOS}{ANEXOS}

\section*{Anexo A. Reproducibilidad experimental}

El repositorio del proyecto contiene el paquete \texttt{topicagg} con sus dependencias
declaradas en \texttt{pyproject.toml} y una suite de 25 pruebas automatizadas; el contrato de
corrida \texttt{cava-run-v1}, con un manifiesto por video (commit de git, indicador de
cambios sin confirmar, perfil, hiperparámetros y \texttt{gt\_used: false}) y 115 corridas
registradas; el registro cronológico de experimentos con sus pre-registros y veredictos
(\texttt{docs/EXPERIMENT\_LOG.md}); la decisión de arquitectura ADR-0001 (sin fusión de
representaciones, votantes activos y abstención); los reportes semanales en
\texttt{reports/}; y los JSON de resultados por experimento en \texttt{results/}.

\section*{Anexo B. Inventario de módulos}

El inventario de los 26 módulos del pipeline (implementación, evidencia y estado), la tabla de
lo que puede afirmarse por tipo de contenido y la lista de conjuntos con anotación disponible
están en \texttt{docs/progress.md}. La Tabla~\ref{tab:inventario} del Capítulo~IV resume ese
inventario.

\section*{Anexo C. Glosario y experimentos citados}

\begin{table}[H]
    \centering
    \caption{Términos y nombres internos usados en el documento.}
    \label{tab:glosario}
    \small
    \begin{tabular}{p{3.2cm}p{10.2cm}}
        \hline
        \textbf{Término} & \textbf{Significado} \\
        \hline
        \texttt{cava\_abst} & Configuración final del gestor: retornos con percentil 90,
                              enlace promedio y abstención. \\
        \texttt{cava}, \texttt{cava\_avg} & Variantes sin abstención, con componentes
                              conexas o con enlace promedio. \\
        \texttt{\_r80}, \texttt{\_r95} & Sufijos que cambian el percentil de retorno a 80 o 95. \\
        \texttt{fused} & Base que concatena los embeddings normalizados de las señales y los
                         trata como una sola señal en el mismo gestor. \\
        \texttt{contiguous} & Gestor sin retornos: solo produce tópicos contiguos. \\
        Proxy de reparto & Referencia de MovieNet: dos tomas son del mismo tópico si
                           comparten personajes. \\
        Procedencia & Aporte de cada señal a un grupo (acuerdo, soporte y coherencia). \\
        Percolación & Colapso del mapa en pocos tópicos gigantes por encadenamiento de
                      enlaces. \\
        P\#, V\# & Identidades por apariencia (no publicadas) y por voz. \\
        blind40, raw\_ext & Subconjuntos de VidChapters para evaluación ciega (40 videos) y
                            desarrollo (50 videos). \\
        S\&W, I\&W & Filtros de habla de EXP-034: intersección de Silero VAD (S) o de ina (I)
                        con la detección de habla de Whisper (W). \\
        Bases U, S, T & Bases de EXP-040: clips equiespaciados, mejor señal sola y
                        segmentación de un modelo de lenguaje sobre los subtítulos. \\
        \hline
    \end{tabular}
\end{table}

\begin{table}[H]
    \centering
    \caption{Experimentos citados en el documento.}
    \label{tab:experimentos}
    \small
    \begin{tabular}{p{2.3cm}p{7.4cm}p{3.6cm}}
        \hline
        \textbf{Experimento} & \textbf{Pregunta} & \textbf{Resultado} \\
        \hline
        EXP-017/018 & ¿El mapa por consenso supera a la fusión y al contiguo en MCC? &
                      Sí con abstención ($z = 2.46$) \\
        EXP-020a & ¿La voz identifica al hablante mejor que la imagen? & Sí (AUC 0.74 frente
                   a 0.55) \\
        EXP-023 & ¿Los turnos de hablante mejoran las fronteras en reuniones? & Parcial ($z =
                  2.32$; empata con contiguo) \\
        EXP-025 & ¿El perfil de documental se acerca al acuerdo humano en BBC? & Sí ($z =
                  6.1$ y $7.4$) \\
        EXP-027 & ¿El ajuste a cortes visuales es preciso? & Sí ($z = 26.6$ y $23.0$) \\
        EXP-028 & ¿Los filtros de habla mejoran la identidad de voz? & Sí (ARI 0.64 y
                  0.73) \\
        EXP-029 & ¿El detector de hablante único evita falsos en reuniones? & Sí (0/40), sin
                  narrador (0/11) \\
        EXP-030 & ¿Parakeet reemplaza a Whisper? & No \\
        EXP-031 & ¿Los destacados siguen la importancia humana en TVSum? & No \\
        EXP-032 & ¿CAVA reproduce los capítulos de autor? & Solo supera a la fusión \\
        EXP-033 & ¿Un nivel de capítulo por fronteras funciona? & Parcial \\
        EXP-037 & ¿Una regla fija de destacados supera a las bases? & No ($z = 0.70$) \\
        EXP-038 & ¿Una política aprendida transfiere entre conjuntos? & No ($z = -2.62$) \\
        EXP-039 & ¿Otro detector resuelve al narrador con habla filtrada? & No ($z = -0.37$) \\
        G-enlace-2 & ¿Un juez reconoce los retornos como el mismo hilo? & No ($p = 0.72$) \\
        EXP-040 & ¿Las personas prefieren las versiones de CAVA? & Diseñado, pendiente \\
        \hline
    \end{tabular}
\end{table}

\section*{Anexo D. Definiciones operativas de métricas}

El MCC del mapa se calcula sobre la clasificación de pares de unidades (mismo tópico o no)
contra la referencia: proxy de reparto en MovieNet y anotación de tópicos en AMI. El margen
de coherencia intra$-$inter compara, por señal, la similitud coseno de cada unidad con su
grupo (dejándola fuera) y la similitud entre grupos. El ARI evalúa la identidad de hablante
contra la anotación de AMI y VoxConverse. La recompensa de clips es el $z$ de la fracción de
anotadores que eligió cada segundo en TVSum y SumMe. G-versión verifica la existencia y el
presupuesto de cada lista por video y formato. Las diferencias se reportan con el error
estándar por video (o por reunión), y el umbral de validación es $z > 2$ salvo que el
pre-registro indique otro.
